\chapter{Conclusion}
\label{ch:conclusion}

\section{Principal Findings}

Three findings stand out as the most broadly informative.

\textbf{Sequence models learn meaningful vulnerability representations, but hit a predictable ceiling.}
The renaming and ablation test results confirm that the model has learned something beyond
keyword association.
But the dead-code and data-flow failures locate the ceiling precisely: anything requiring
inter-statement reasoning lies outside what a sequence encoder can represent.
This ceiling is an inherent property of the architecture, not a fixable training artefact.
Token attribution analysis makes the mechanism explicit: CWE classes with locally visible
signals produce correct attributions, while data-flow-dependent classes produce diffuse or
surface-anchored attributions that explain exactly why the architectural failures occur.

\textbf{Cross-language transfer quality is determined by surface representation, not semantic equivalence.}
CWE-798 transfers perfectly in both directions because a string literal assigned to a
password variable looks identical in C and Python.
CWE-367 transfers almost not at all because the C and Python realisations use structurally
different API sequences.
This clarifies what ``cross-language vulnerability detection'' means in practice:
detecting shared surface patterns, not shared semantics.

\textbf{Joint training consistently outperforms single-language training.}
Across all three test sets the joint model outperforms both single-language variants.
The reason is mutual reinforcement: the eight CWE classes shared between C and Python
provide the model with two structurally distinct surface realisations of the same
vulnerability concept during training, which strengthens the underlying representation
rather than diluting it.
Single-language models over-specialise --- the Python-only model achieves F1\,=\,0.976
on Python but collapses to 0.677 on C, having never seen C memory-management idioms.
The joint model avoids this by using each language's examples as implicit regularisation
for the other, producing representations that generalise better even within each language
individually: the joint model's C F1 (0.905) exceeds the C-only model's (0.861) because
Python examples for shared CWEs reinforce the corresponding C patterns during training.

\textbf{Real-world augmentation is the highest-leverage single intervention.}
The 0.278 F1 point improvement from adding 1{,}000 real-world samples dwarfs all other
single changes.
For practitioners building on similar approaches, the implication is direct: if real-world
performance matters, investing in labelled real-world training data has a higher return
than architectural experimentation on synthetic benchmarks alone.

\section{Closing Remarks}

The system demonstrates that a pre-trained code model, with careful dataset design,
targeted hard negatives, and real-world augmentation, can reach practically useful
performance on a multi-language vulnerability detection task.
The semantic understanding tests move the evaluation beyond accuracy figures to a
qualitative characterisation of what the model has actually learned, which is a
precondition for trusting any machine learning system in a security-critical context.

\textbf{Limitations on complex vulnerability detection.}
The results make clear that the model handles locally visible patterns reliably but
struggles as vulnerability complexity increases.
CWE classes whose exploitability depends on a single self-contained expression ---
hardcoded credentials, weak algorithm selection, integer overflow in an arithmetic
expression --- reach perfect or near-perfect F1.
Classes that require reasoning across multiple statements degrade substantially: CWE-416
(use-after-free) requires tracking a pointer from its allocation through its
deallocation to a subsequent dereference; CWE-367 (TOCTOU) requires recognising a
temporal gap between a check and its use.
The dead-code test failures confirm this directly: the model cannot determine whether a
dangerous code path is reachable, so it flags unreachable vulnerable blocks as if they
were live.
Extending this approach to more complex vulnerability classes --- inter-procedural
taint tracking, concurrency bugs, logic flaws spanning multiple files --- would
require architectural changes that bring explicit program structure into the
model's representational vocabulary, such as graph-augmented encoding over data-flow
and control-flow graphs.

\textbf{Positioning within the broader research landscape.}
This work sits within a rapidly evolving field.
Recent multilingual benchmarking by Shu et al.~\cite{shu2026multilingual} confirms that
cross-language vulnerability detection remains unsolved even for the largest pre-trained
models, with accuracy dropping substantially outside the C/C++ domain.
Nguyen et al.~\cite{nguyen2025mulvuln} similarly demonstrate that language-specific
nuances persist as a meaningful gap for multilingual transformer models.
The contributions here --- a focused two-language fine-tuning study, a diagnostic
semantic test suite, and token-level attribution analysis --- complement that body of
work by providing a principled characterisation of \emph{why} sequence models fail on
specific classes, rather than only reporting where they fail.
As the field moves towards graph-augmented architectures, LLM-assisted code repair, and
automated vulnerability patching, the baseline and diagnostic framework established
here provide a reference point for measuring genuine progress: a new architecture that
solves the dead-code and data-flow cases identified in this work would represent a
qualitative advance, not merely a benchmark improvement.

The gap between synthetic and real-world performance is not a failure of the approach
but a measurement of the problem's difficulty.
Closing it further requires more real-world training data, broader CWE coverage, and ---
ultimately --- architectural advances that bring data-flow and control-flow reasoning
into the model's representational vocabulary.
The results reported here establish a clear baseline, locate the ceiling of sequence-based
approaches precisely, and map a concrete path forward.

\section{Future Directions}

Five directions would most directly address the identified limitations and extend the
scope of this work.

\textbf{Graph-augmented encoding.}
Replacing or augmenting the sequence encoder with a graph neural network operating on the
program's data-flow graph would enable the model to track taint directly --- following
user-controlled values from input sources through assignments and function calls to
execution sinks.
This would be the most impactful single architectural change, directly addressing the
CWE-78 data-flow failure and likely improving CWE-416 (use-after-free) and CWE-476
(null dereference).

\textbf{Multi-class CWE classification.}
The current system is a binary classifier that predicts whether a function is vulnerable
or safe.
A natural extension is a multi-class classifier that simultaneously predicts which CWE
class a vulnerable function belongs to.
This would be considerably more useful in practice: a tool that reports ``this function
contains a use-after-free (CWE-416)'' provides immediately actionable information,
whereas a binary VULNERABLE verdict requires a follow-up manual triage step.
The fine-tuned model's CLS-token embeddings already encode CWE-discriminative
information --- the per-CWE F1 variation across classes demonstrates this --- and a
lightweight multi-class head trained on top of those embeddings is a direct next step.
The main challenge is label quality: real-world datasets such as DiverseVul provide
CWE annotations for only a subset of samples, and synthetic data would need to be
extended to cover a wider range of classes at sufficient volume per class to train a
reliable multi-class head.

\textbf{Larger real-world training corpus.}
Using more of the DiverseVul training split --- or combining DiverseVul with BigVul and
additional CVE-labelled repositories --- would reduce the generalisation gap further.
The 0.278 F1 improvement from 1{,}000 real samples suggests a favourable response to
real data volume; a 10{,}000-sample augmentation might plausibly close the remaining gap
substantially.

\textbf{CWE class expansion.}
Adding synthetic training examples for CWE-787 (out-of-bounds write), CWE-400 (uncontrolled
resource consumption), and CWE-703 (improper exception handling) would cover the three most
common false-negative classes in the real-world test.
Template-based synthetic generation for these classes follows the same methodology used
for the existing 17 CWEs.

\textbf{Multi-language real-world evaluation.}
Constructing or obtaining a Python-language real-world vulnerability evaluation set ---
sourced from PyPI security advisories or OSV database entries --- would allow the same
synthetic-to-real gap analysis for Python that Experiment~3 provides for C.
